% --- Conteúdo para conclusao.tex ---

\qquad Este trabalho analisou cinco algoritmos de ordenação sob três cenários de entrada, confrontando a análise teórica de complexidade com os tempos obtidos experimentalmente. Os resultados confirmaram o comportamento previsto pela teoria e evidenciaram o impacto prático da escolha do algoritmo em aplicações que lidam com grandes volumes de dados.

\qquad A principal constatação diz respeito à separação entre os algoritmos quadráticos e os mais eficientes, \textit{Shell Sort} e \textit{Merge Sort}. Enquanto \textit{Bubble Sort}, \textit{Insertion Sort} e \textit{Selection Sort} levaram entre 18 e 66 minutos para ordenar um milhão de elementos aleatórios, o \textit{Shell Sort} concluiu a mesma tarefa em 0,461 segundo e o \textit{Merge Sort} em 0,292 segundo. Essa diferença não decorre de otimizações no código, e sim da ordem de complexidade de cada método: o ganho obtido ao substituir um algoritmo $O(n^2)$ por um subquadrático ou $O(n \log n)$ supera, em larga escala, qualquer ajuste pontual de implementação.

\qquad O \textit{Merge Sort} foi o único algoritmo recursivo do trabalho e o único baseado na estratégia de dividir para conquistar. A resolução de sua recorrência resultou em $T(n) = n + n \log_2 n$, ou seja, $O(n \log n)$ nos três casos, e os tempos medidos confirmaram essa previsão: ao aumentar a entrada de 100.000 para 1.000.000 de elementos, o tempo do caso aleatório cresceu cerca de 12,7 vezes, valor próximo das 12 vezes esperadas para $n \log n$. Também ficou evidente que sua eficiência tem um custo, pois ele exige um vetor auxiliar de tamanho $n$, enquanto os demais algoritmos ordenam o vetor \textit{in-place}.

\qquad A comparação entre os dois algoritmos mais rápidos mostrou que não existe um vencedor absoluto. O \textit{Merge Sort} foi o mais rápido com entrada aleatória e apresentou tempos parecidos em todos os cenários, já que realiza o mesmo trabalho em qualquer ordem dos dados. O \textit{Shell Sort}, por outro lado, foi mais rápido nas entradas crescente e decrescente, pois aproveita a organização prévia do vetor, e não utiliza memória adicional. A escolha entre eles depende, portanto, do tipo de dado esperado e da memória disponível.

\qquad Os testes também mostraram que a ordem inicial dos dados afeta cada algoritmo quadrático de maneira distinta. O \textit{Insertion Sort} foi o mais sensível, variando de 0,005 segundo no melhor caso a 2.164,840 segundos no pior, diferença de mais de quatrocentas mil vezes, o que o torna vantajoso quando há expectativa de dados parcialmente ordenados. O \textit{Selection Sort} apresentou o comportamento mais previsível entre os quadráticos, com variação de apenas 21\% entre os cenários. Já o \textit{Bubble Sort} não obteve vantagem em nenhum cenário, tendo registrado o pior desempenho nas entradas decrescente e aleatória.

\qquad Um resultado que merece registro é o fato de o \textit{Bubble Sort} ter sido ligeiramente mais lento com entrada aleatória do que com entrada decrescente, contrariando à primeira vista a expectativa teórica de pior caso. A análise dessa ocorrência mostrou que, na implementação utilizada, o número de comparações é idêntico nos três cenários, variando apenas o número de trocas e o padrão de acesso à memória. Isso reforça que a análise assintótica descreve a tendência de crescimento do algoritmo, e não o tempo exato de execução, que depende também de fatores de hardware.

\qquad Observou-se ainda que, para entradas de até 1.000 elementos, os cinco algoritmos concluíram a ordenação em tempo inferior a 0,005 segundo, faixa em que a escolha entre eles é irrelevante. As diferenças só passam a ser determinantes a partir de dezenas de milhares de elementos, o que indica que a seleção de um algoritmo deve considerar a escala esperada do problema, e não apenas sua classificação teórica.

\qquad Por fim, o trabalho permitiu articular os conteúdos de análise de complexidade com a prática de implementação e medição. A construção das tabelas de custo, dos testes de mesa, da recorrência do \textit{Merge Sort} e das funções de tempo para cada cenário tornou possível prever o comportamento dos algoritmos antes mesmo da execução, e os resultados experimentais confirmaram essas previsões. Como continuidade natural do estudo, seria pertinente incluir o \textit{Quick Sort}, avaliar versões otimizadas do \textit{Merge Sort}, como o uso do \textit{Insertion Sort} em trechos pequenos, e testar o efeito de diferentes sequências de \textit{gaps} sobre o desempenho do \textit{Shell Sort}.